% !TeX program = xelatex
\documentclass[UTF8,a4paper,zihao=-4,fontset=windows]{ctexart}
\usepackage[margin=23mm]{geometry}
\usepackage{booktabs,tabularx,hyperref,listings,xcolor}
\hypersetup{colorlinks=true,linkcolor=black,urlcolor=black}
% Cover metadata supplied by the user.
\newcommand{\MYTITLE}{作业1}
\newcommand{\MYTERM}{2026--2027 学年第一学期}
\newcommand{\MYCOURSE}{专业综合实践}
\newcommand{\MYCOURSEID}{0610140}
\newcommand{\MYADVISOR}{王卯宁}
\newcommand{\MYCLASS}{}
\newcommand{\MYID}{}
\newcommand{\MYNAME}{}
% Full teaching-class identifier: (2026-2027-1)-0610140-01

\begin{document}
% Cover adapted from Xueyan Hu's CUFE Course Report template (CC BY 4.0).
% https://www.overleaf.com/latex/templates/central-university-of-finance-and-economics-cufe-course-report-latex-template/zjsmnqbwgmyg
\newgeometry{a4paper,left=2.5cm,right=2.5cm,top=2.54cm,bottom=2.54cm}
\begin{titlepage}
\thispagestyle{empty}
\begin{spacing}{1.5}
\fontsize{12pt}{14.5pt}\selectfont
\begin{center}
    \includegraphics[width=1.0\textwidth]{template/Im0.png}\\
    \vspace{20mm}
    \vspace{10mm}
    \setlength{\extrarowheight}{3mm}
    {\kaishu\zihao{4}
    \begin{tabular}{rp{8cm}<{\centering}}
        {\makebox[4\ccwd][s]{学年学期：}} & \kaishu\underline{\makebox[8cm]{\MYTERM}} \\
        {\makebox[4\ccwd][s]{课程名称：}} & \kaishu\underline{\makebox[8cm]{\MYCOURSE}} \\
        {\makebox[4\ccwd][s]{课程代码：}} & \kaishu\underline{\makebox[8cm]{\MYCOURSEID}} \\
        {\makebox[4\ccwd][s]{任课教师：}} & \kaishu\underline{\makebox[8cm]{\MYADVISOR}} \\
        {\makebox[4\ccwd][s]{班\qquad 级：}} & \kaishu\underline{\makebox[8cm]{\MYCLASS}} \\
        {\makebox[4\ccwd][s]{学\qquad 号：}} & \kaishu\underline{\makebox[8cm]{\MYID}} \\
        {\makebox[4\ccwd][s]{姓\qquad 名：}} & \kaishu\underline{\makebox[8cm]{\MYNAME}} \\
        \\
        {\makebox[4\ccwd][s]{总\qquad 分：}} & \kaishu\underline{\makebox[8cm]{}} \\
        {\makebox[4\ccwd][s]{评$\ $分$\ $人：}} & \kaishu\underline{\makebox[8cm]{}} \\
    \end{tabular}
    }\\[2cm]
\end{center}
\end{spacing}
\end{titlepage}
\restoregeometry
\setcounter{page}{1}

\begin{center}
{\Large\bfseries 作业1：TuGraph 电影知识图谱实验报告}
\end{center}
\section{实验目标}
使用 TuGraph 4.5.1 建立电影知识图谱，掌握 CSV 属性图建模、数据导入、Cypher 增删改查和聚合分析。采用 Movie、Person 两类节点，通过 PRODUCED 连接制作人员与电影。
\section{实验数据}
数据来自三份电影业务 CSV：movie.csv、person.csv、produce.csv。原始文件分别含有 54 部电影、665 名人物和 60 条制作关系（文件记录数量，不是数据库实测数量）。电影主键为 id，属性包括 title、tagline、summary、poster\_image、duration 和 rated；人物主键为 id，属性包括 name、born、poster\_image。PRODUCED 由 Person 指向 Movie，端点以主键匹配。
\section{实验环境与流程}
在 Windows PowerShell 和 Docker Desktop 的 Linux x86\_64 容器中使用 TuGraph 4.5.1。先运行 \texttt{scripts/prepare\_movie\_csv.py}，将源 CSV 转为统一编码、标准表头的文件。导入配置位于 \texttt{data/movie\_csv/import.json}。执行 \texttt{setup.ps1} 完成镜像构建与离线导入，图名为 \texttt{movie\_csv\_demo}，数据库保存在 \texttt{runtime/db}。
\section{查询设计}
\subsection{基本数量与关联查询}
通过 \texttt{MATCH (m:Movie) RETURN count(m)} 统计电影节点数量；通过 \texttt{MATCH (p:Person)-[:PRODUCED]->(m:Movie)} 查询人物与电影制作关系。
\subsection{制作人员作品数量}
\begin{lstlisting}[basicstyle=\ttfamily\small,breaklines=true]
MATCH (p:Person)-[:PRODUCED]->(m:Movie)
RETURN p.name AS producer, count(DISTINCT m) AS movie_count
ORDER BY movie_count DESC, producer LIMIT 20;
\end{lstlisting}
\subsection{电影制作人员数量}
\begin{lstlisting}[basicstyle=\ttfamily\small,breaklines=true]
MATCH (p:Person)-[:PRODUCED]->(m:Movie)
RETURN m.title AS movie, count(DISTINCT p) AS producer_count
ORDER BY producer_count DESC, movie LIMIT 20;
\end{lstlisting}
\subsection{影片时长与分级}
利用 \texttt{avg(m.duration)} 和 \texttt{count(m)} 分析不同 \texttt{rated} 分组的片长与电影数量；分析仅针对数据中具有时长的电影。
\section{增删改查}
建立主键为 -999999 的临时 Movie 节点，再读取、更新其 title、删除并检查剩余数量。该临时记录不属于原始电影数据，在操作结束后应为 0。
\section{复现与结果核对}
在 Browser 运行 \texttt{cypher/movie\_csv.cypher} 中的语句，并执行 \texttt{python scripts/verify\_movie\_csv.py}。脚本自动比较导入 CSV 的行数和实时数据库查询结果，保存 JSON 供实验核查。
\textbf{当前版本未附数据库运行截图或实测输出。} 正式提交前应在实际部署环境运行验证，并在此处补充结果及截图；不得将 CSV 行数当作数据库实际验证结果。
\section{实验总结}
本实验以电影及制作人员的属性图表示实体与制作关系，并将 CSV 数据准备、TuGraph 导入、关系查询、聚合和 CRUD 组织为可复现流程。数量一致性及运行截图需以实际部署为依据。
\end{document}
